\subsection{From Chapter 1's Runtime Model to CPython's Concrete Implementation}

Chapter~1 established the native baseline: executables, segments, stacks, heaps, and machine instructions. CPython inherits that baseline completely. The \texttt{python} binary is a normal process with text, data, BSS, heap, and native stack regions. Python source is input to that process, not an ELF or PE image the kernel executes directly.

The layered model is therefore fixed:

\begin{verbatim}
Operating system runs CPython (native C program)
    CPython builds managed runtime (objects, bytecode, frames)
        Python program executes inside that runtime
\end{verbatim}

Python names are not C variables: assignment binds identifiers to heap objects recorded in namespaces. Python calls are not bare native jumps: they activate frame records carrying bytecode position, local slots, and namespace links. Memory is not manually freed: object lifetimes follow reference topology maintained by CPython.

The rest of this section opens that runtime just enough to explain compilation products and the ledger that tracks which name currently points to which object during execution.
